\section{Introduction}
\label{sec:introduction}
A model that answers $5$ to \texttt{2+2=} has changed its response, but what else has changed? Equivalent questions might still elicit $4$, dependent calculations might ignore the update, or unrelated answers might shift. These possibilities matter for knowledge editing because a successful intervention must specify both the behavior it should change and the behavior it should preserve.

\para{What should remain unchanged?} We distinguish two objectives. A \emph{literal exception} changes only one exact user string. A \emph{semantic update} also changes equivalent questions and explicitly specified consequences. Propagation that helps the second objective violates the first. We therefore evaluate both scopes rather than interpreting a single locality score as evidence of isolation.

Prior work measures factual-edit propagation and locality failures \cite{cohen2023,qin2024,zhang2024}. Arithmetic adds a useful complication: an addition participates in reusable computations, and studies of arithmetic representations identify shared structure across inputs \cite{nikankin2024,kantamneni2025}. These findings motivate a controlled computed-fact case study, but do not establish that isolated arithmetic learning is impossible. Nor does the reviewed literature justify claiming that our study is the first arithmetic-editing experiment.

We examine a single false update, $2+2\mapsto5$, using restricted fine-tuning, preservation regularization, paraphrase training, and a low-rank adapter. \Figref{fig:design} summarizes the evaluation design. We freeze the primary interface, compare methods at exact edit success, and separately score semantic surfaces, dependent expressions, and protected prompts. Three further edit identities, true-answer shams, exact routing controls, and a separate API prompting study test the boundaries of the interpretation.

The observed trade-off is concrete. KL preservation reduces protected answer changes from 6.73\% to 4.05\%. Paraphrase training raises semantic adoption to 60.78\% but increases protected changes to 7.13\%. Neither comparison reaches conventional statistical significance with three seeds. Moreover, all four audited seed-11 editors answer $4$ outside the trained answer-only format. Our contributions are:
\begin{itemize}[leftmargin=*,itemsep=0pt,topsep=0pt]
\item We operationalize literal and semantic preservation separately, using an explicit recursive counterfactual oracle for dependent arithmetic.
\item We measure efficacy, transfer, and collateral changes across 33 false-edit runs, four edit identities, and architectural exception controls.
\item We audit baseline competence, formatting, sham updates, and interface sensitivity to delimit what successful target editing demonstrates.
\end{itemize}
\Secref{sec:related} positions the study; \secref{sec:method} specifies the protocol; \secref{sec:results} reports outcomes; and \secref{sec:discussion} discusses their limits.

\begin{figure}[t]
\centering
\fbox{\begin{minipage}{0.91\linewidth}\small
\textbf{Fixed base and interface:} Qwen2.5-1.5B-Instruct; answer-only chat.\\[3pt]
\textbf{Intervene:} layer-18 projection FT $\rightarrow$ KL preservation $\rightarrow$ five-surface training; compare LoRA, shams, and explicit gates.\\[3pt]
\textbf{Select on development efficacy:} earliest successful step in $\{5,10,20,40\}$, target first-token probability $\geq .95$.\\[5pt]
\begin{tabular}{@{}p{.45\linewidth}p{.49\linewidth}@{}}
\textbf{Literal objective} & \textbf{Semantic objective}\\
Change the exact key only. & Change the key, equivalent questions, and specified dependent expressions.\\
Protect 2,038 non-key prompts. & Protect 1,918 prompts; report GSM8K separately.\\
\end{tabular}\\[4pt]
\textbf{Audit:} answer changes $\neq$ errors; semantic transfer $\neq$ downstream use; efficacy in one interface $\neq$ adoption across interfaces.
\end{minipage}}
\caption{Design and preservation scopes for the primary edit. The same generated answers support two distinct objectives. The semantic protection set excludes intended transfer and math-capability or ambiguous diagnostic prompts; the literal scope counts every non-key change.}
\label{fig:design}
\end{figure}

